% === Begin Label Data ===
% ID: 238
% 难度星级: 
% 标签: 
% 备注: 
% 组卷引用次数: 0
% === End  Label Data ===

\begin{problem}{2011}{G}{安徽卷（理）}{18}{数列，三角函数}
在数 $1$ 和 $100$ 之间插入 $n$ 个实数, 使得这 $n+2$ 个数构成递增的等比数列, 将这 $n+2$ 个数的乘积记作 $T_n$, 再令 $a_n = \lg T_n, n \ge 1$.

（1）求数列 $\{a_n\}$ 的通项公式;

（2）设 $b_n = \tan a_n \cdot \tan a_{n+1}$, 求数列 $\{b_n\}$ 的前 $n$ 项和 $S_n$.
\end{problem}

\begin{solutions}
(1) 设 $t_1, t_2, \cdots, t_{n+2}$ 构成等比数列，其中 $t_1=1, t_{n+2}=100$，则
$$
T_n=t_1 t_2 \cdots t_{n+1} t_{n+2}, T_n=t_{n+2} t_{n+1} \cdots t_2 t_1.
$$

两式相乘，并利用 $t_i t_{n+3-i} = t_1 t_{n+2} = 10^2 (1 \leqslant i \leqslant n+2)$，可得
$$
T_n^2=(t_1 t_{n+2})(t_2 t_{n+1}) \cdots (t_{n+1} t_2)(t_{n+2} t_1) = 10^{2(n+2)}.
$$

从而
$$
a_n=\lg T_n=n+2 (n \geqslant 1).
$$

(2) 依题意并结合(1)中的计算结果，可知
$$
b_n=\tan(n+2) \cdot \tan(n+3) (n \geqslant 1).
$$

利用 $\tan 1 = \tan[(k+1)-k] = \dfrac{\tan(k+1)-\tan k}{1+\tan(k+1) \cdot \tan k}$，得
$$
\tan(k+1) \cdot \tan k = \frac{\tan(k+1)-\tan k}{\tan 1} - 1.
$$

所以
$$
S_n = \sum_{k=1}^n b_k = \sum_{k=3}^{n+2} \tan(k+1) \cdot \tan k = \sum_{k=3}^{n+2} \left[\frac{\tan(k+1)-\tan k}{\tan 1} - 1\right] = \frac{\tan(n+3)-\tan 3}{\tan 1} - n.
$$
\end{solutions}